\documentclass[11pt]{article}
% Uses the canonical latex-docs semantic style in-repository. The local
% lightweight fallback is for independent compilation of this distribution.
\IfFileExists{technical-implementation.sty}{\usepackage{technical-implementation}}{\usepackage{../support/aaoa-standalone}}
\usepackage{mathtools}
\setDocTitle{Python Laboratories}
\setDocSubtitle{Four computational experiments and reproducible evidence}
\setDocVersion{1.0}
\setDocStatus{Educational reference / unapproved working material}
\setDocOwner{Algorithm Engineering}
\setDocAuthor{Original companion material}
\setDocSummary{Standalone Python laboratory instructions for sorting, generating functions, trees, tries and occupancy.}
\providecommand{\Expect}{\mathbb{E}}
\providecommand{\Prob}{\mathbb{P}}
\providecommand{\Var}{\operatorname{Var}}
\providecommand{\bigoh}{\mathcal{O}}
\begin{document}
\makedoctocpage

\section{Toolchain and lab contract}
Python 3.10 or newer, standard library only for core laboratories. Optional SymPy is recommended for formal-series checks but is not required by the supplied scripts. Always save source revision and interpreter version, seed, generated data description, operation counts and raw benchmark observations. Do not time input construction.
\section{LAB-PY01: average-case quicksort}
\textbf{Hypothesis:} under a uniformly random pivot rank, the expected number of key comparisons for the supplied partition model is $2(n+1)H_n-4n$ (for $n\ge1$). Run all distinct permutations of sizes up to six using a fixed-first-pivot algorithm; compare the enumeration mean to the formula. Then draw seeded random permutations for larger $n$.
\textbf{Procedure:} (1) instrument comparisons, (2) validate sorted output, (3) enumerate input permutations, (4) report mean, variance and at least one adversarial case, (5) explain where $n$ versus number of \emph{distinct} keys matters. A timing result is secondary to the exact comparison count.
\section{LAB-PY02: generating-function and asymptotic checks}
\textbf{Hypothesis:} Catalan coefficients satisfy\[C_0=1,\qquad C_{n+1}=\sum_{k=0}^n C_kC_{n-k},\qquad \frac{C_n\sqrt{\pi}n^{3/2}}{4^n}\to1.\]Calculate exact values up to $n=50$ using integer arithmetic; compare to the asymptotic expression using logarithms to avoid overflow. As an independent check, verify $[z^n]z/(1-z)^3=n(n+1)/2$ through order eight.
\section{LAB-PY03: random BSTs are not uniform Catalan trees}
\textbf{Hypothesis:} random insertion permutations produce expected internal path length $2(n+1)H_n-4n$, but they do not assign equal probability to every Catalan shape. Enumerate all permutations for $n\le6$, record shape frequencies, calculate path length and compare to a uniform average over distinct shapes. A tree's shape is a nested pair; it is not a set of keys.
\section{LAB-PY04: tries and occupancy}
\textbf{Hypothesis:} longer shared prefixes reduce the number of newly allocated nodes in a basic character trie. Construct (A) disjoint-prefix keys and (B) common-prefix keys of equal count/length; record node count and characters inspected per lookup. Separately sample independent uniform bucket placements; compare mean occupancy $m/b$ and pair-collision estimate $\binom m2/b$.
\section{Commands and validation}
From \texttt{laboratories/code}, run the Python tests with \texttt{python3 -m unittest -v test\_models.py}. For a complete benchmark command and CSV options, consult \texttt{code/README.md}. No package installation is required. CSV is generated measurement data, not a source file to commit without a purpose.
\section{Required report}
For each lab include: model and variables; formal prediction; test oracle; program revision; observed data; normalized cost; uncertainty when estimated; adversarial or alternate-model check; limitation; and recommendation. Use the laboratory-report template and experiment-plan template.

\end{document}
